\chapter{Elliptic estimates for weak solutions}\label{sec:elliptic-tools}
The estimates below separate three levels of information: energy control, continuity with merely measurable coefficients, and higher regularity with continuous coefficients. They supply the analytic inputs used in the minimal graph and Jacobi-field sections. All balls in this section are Euclidean balls in $\mathbb R^n$, with $n\geq2$.

\begin{definition}[Sobolev spaces and weak derivatives]\label{def:sobolev-weak-derivatives}
For an open $U\subset\mathbb R^n$, $1\leq p<\infty$, and integer $k\geq0$, the space $W^{k,p}(U)$ consists of functions whose distributional derivatives $D^\beta u$, $|\beta|\leq k$, belong to $L^p(U)$. Its norm is $\sum_{|\beta|\leq k}\|D^\beta u\|_{L^p(U)}$. A first weak derivative is characterized by $\int uD_i\phi=-\int D_i u\,\phi$ for all $\phi\in C_c^\infty(U)$. We write $H^1=W^{1,2}$, and $W^{1,p}_0(U)$ for the closure of $C_c^\infty(U)$ in $W^{1,p}(U)$. The subscript $\mathrm{loc}$ requires the corresponding property on every relatively compact subdomain. See \cite{SimonPDE}, Chapter 5.
\end{definition}

\begin{definition}[Uniform ellipticity and weak divergence equations]\label{def:weak-elliptic-equation}
Let $A:U\to\mathbb R^{n\times n}$ be measurable and symmetric. We assume that constants $0<\lambda\leq\Lambda<\infty$ satisfy
\[
 \lambda|\xi|^2\leq A(x)\xi\cdot\xi\leq\Lambda|\xi|^2
 \quad\text{for almost every }x\in U,\ \xi\in\mathbb R^n.
\]
A function $v\in H^1_{\mathrm{loc}}(U)$ solves $\operatorname{div}(A Dv)=0$ weakly if $\int A Dv\cdot D\phi=0$ for all $\phi\in C_c^\infty(U)$. A weak subsolution of this equation satisfies $\int A Dv\cdot D\phi\leq0$ for nonnegative test functions. No derivatives or continuity of $A$ are assumed. This is the homogeneous, symmetric, zero-lower-order-term case of \cite{SimonPDE}, Chapters 6 and 17.
\end{definition}

\begin{definition}[Hölder seminorms]\label{def:holder-seminorm}
For $0<\alpha<1$, define
\[
 [f]_{\alpha;U}=\sup_{x,y\in U,\ x\ne y}\frac{|f(x)-f(y)|}{|x-y|^\alpha}.
\]
The space $C^{k,\alpha}(U)$ consists of $C^k$ functions with finite norm $\sum_{|\beta|\leq k}\|D^\beta f\|_{C^0(U)}+\sum_{|\beta|=k}[D^\beta f]_{\alpha;U}$. Local versions are defined on relatively compact subdomains. On a ball of radius $R$, the scale-invariant combination for the $k$th derivative is $R^k\|D^kf\|_\infty+R^{k+\alpha}[D^kf]_\alpha$. See \cite{SimonPDE}, Chapter 12.
\end{definition}

\begin{lemma}[Caccioppoli energy estimate]\label{lem:caccioppoli}
If $v$ solves the equation in Definition \ref{def:weak-elliptic-equation}, $\eta\in C_c^1(U)$, and $c\in\mathbb R$, then
\[
 \int\eta^2|Dv|^2\leq4(\Lambda/\lambda)^2\int|D\eta|^2|v-c|^2.
\]
In particular, if $B_R\subset U$ and $0<r<R$, then $\int_{B_r}|Dv|^2\leq C(R-r)^{-2}\int_{B_R}|v-c|^2$, with $C$ depending only on $\Lambda/\lambda$ and the cutoff choice.
\end{lemma}
\begin{proof}
By density we may test with $\eta^2(v-c)$. Ellipticity and the operator-norm bound $|A|\leq\Lambda$ give
\[
 \lambda\int\eta^2|Dv|^2
 \leq 2\Lambda\int |\eta||v-c||Dv||D\eta|
 \leq\frac\lambda2\int\eta^2|Dv|^2+
 \frac{2\Lambda^2}{\lambda}\int|v-c|^2|D\eta|^2.
\]
Absorb the first term. A cutoff equal to one on $B_r$ and with $|D\eta|\leq C/(R-r)$ gives the last assertion. Compare the energy-test method in \cite{SimonPDE}, Chapter 17.
\end{proof}

\begin{theorem}[Local boundedness of weak solutions]\label{thm:elliptic-local-boundedness}
Let $v\in H^1(B_{2R}(a))$ solve the equation in Definition \ref{def:weak-elliptic-equation}. Then
\[
 \operatorname*{ess\,sup}_{B_R(a)}|v|
 \leq C\left(R^{-n}\int_{B_{2R}(a)}|v|^2\right)^{1/2},
\]
where $C=C(n,\Lambda/\lambda)$. This is \cite{SimonPDE}, Chapter 18, Theorem 2, p.~212, with zero forcing and lower-order terms. The local iteration proof is supplied below.
\end{theorem}
\begin{proof}
By Lemma \ref{lem:positive-part-subsolution}, $v_+$ and $(-v)_+$ are nonnegative weak subsolutions of the same zero-lower-order equation. Apply Theorem \ref{thm:moser-subsolution-bound} to each. Their $L^2$ norms are at most that of $v$, and $|v|$ is their pointwise maximum. This gives the displayed estimate. The supporting argument uses truncated power tests, the external Sobolev gain, an integrability-improvement step and a convergent product of iteration constants; see Lemmas \ref{lem:truncated-power-energy} and \ref{lem:moser-integrability-step}.
\end{proof}

\begin{theorem}[Interior Harnack inequality]\label{thm:elliptic-harnack}
If $v\in H^1(B_{2R}(a))$ is nonnegative and solves the equation in Definition \ref{def:weak-elliptic-equation}, then
\[
 \operatorname*{ess\,sup}_{B_R(a)}v\leq
 C\operatorname*{ess\,inf}_{B_R(a)}v,
 \qquad C=C(n,\Lambda/\lambda).
\]
The sign assumption is essential, and the estimate is on a smaller ball. This is the zero-lower-order-term specialization of \cite{SimonPDE}, Chapter 18, Theorem 1, p.~211. The proof remains external.
\end{theorem}

\begin{corollary}[Hölder continuity from oscillation decay]\label{cor:elliptic-holder}
There are $\alpha=\alpha(n,\Lambda/\lambda)\in(0,1)$ and $C=C(n,\Lambda/\lambda)$ such that every weak solution on $B_{2R}(a)$ has a $C^{0,\alpha}$ representative on $B_R(a)$ and
\[
 \|v\|_{C^0(B_R(a))}+R^\alpha[v]_{\alpha;B_R(a)}
 \leq C\left(R^{-n}\int_{B_{2R}(a)}|v|^2\right)^{1/2}.
\]
\end{corollary}
\begin{proof}
Local boundedness first bounds $v$ on $B_{3R/2}(a)$. On any ball $B_{2s}$ compactly contained there, let $M$ and $m$ be its essential supremum and infimum. Apply Harnack to $M-v$ and $v-m$. Adding the resulting inequalities gives
\[
 \operatorname*{ess\,osc}_{B_s}v\leq
 \frac{C_H-1}{C_H+1}\operatorname*{ess\,osc}_{B_{2s}}v,
\]
where we may take $C_H>1$. Iteration on concentric dyadic balls gives a bound $C(s/R)^\alpha$ for the oscillation, choosing a smaller exponent if necessary to have $\alpha<1$. The essential oscillations therefore define a unique continuous representative. Balls centered in $B_R(a)$ of radius at most $R/4$, followed by the bound for pairs farther apart, give the stated seminorm estimate. This is the oscillation mechanism of De Giorgi--Nash--Moser theory; see \cite{SimonPDE}, Chapters 17--18.
\end{proof}

\begin{theorem}[Interior Schauder regularity]\label{thm:interior-schauder}
Let $0<\alpha<1$, and let $a^{ij}=a^{ji}$, $b^i$, and $c$ belong to $C^{0,\alpha}(B_1)$, with $\lambda I\leq(a^{ij})\leq\Lambda I$ and the sum of their $C^{0,\alpha}$ norms at most $K$. If $u\in C^2(B_1)\cap C^0(\overline{B_1})$ solves
\[
 a^{ij}D_{ij}u+b^iD_i u+cu=f,\qquad f\in C^{0,\alpha}(B_1),
\]
then $u\in C^{2,\alpha}(B_{1/2})$ and
\[
 \|u\|_{C^{2,\alpha}(B_{1/2})}\leq C\bigl(\|u\|_{C^0(B_1)}+\|f\|_{C^{0,\alpha}(B_1)}\bigr),
\]
with $C=C(n,\alpha,\lambda,\Lambda,K)$. The corresponding higher interior estimates hold when the coefficients and forcing have higher Hölder regularity. No sign assumption on $c$ is needed for this interior estimate. Sources: \cite{SimonPDE}, Chapter 12; \cite{ChMin}, Appendix A, equation (A.2), printed p.~86. The proof remains external.
\end{theorem}

\begin{theorem}[Interior Sobolev estimate for the Laplacian]\label{thm:laplace-w2p}
Let $1<p<\infty$ and $u\in W^{2,p}(B_1)$. Then
\[
 \|u\|_{W^{2,p}(B_{1/2})}\leq C(n,p)
 \bigl(\|u\|_{L^p(B_1)}+\|\Delta u\|_{L^p(B_1)}\bigr).
\]
This is the $k=0$ case of \cite{ChMin}, Appendix A, equation (A.1), printed p.~86. It is an a priori estimate in the stated Sobolev class; neither endpoint $p=1,\infty$ is included. The proof remains external.
\end{theorem}
